\section{Removing a redundant Pell factor}
\label{nine-section}

The central improvement replaces a square of an already certified square
by that square itself.  The square root needed to explain the replacement
exists only in the proof; it does not become an additional unknown.
This changes the arithmetic certificate, so it requires a new soundness
argument.  It produces two complete division-free polynomial circuits:
\begin{center}
\begin{tabular}{lrrr}
Construction & Positive unknowns & Exact degree & Operations\\\hline
Frozen H2082 shared-product parent &9&1,845,940&440\\
Saturated certificate, simple bound &9&1,628,768&434\\
Saturated certificate, whole radical &9&1,103,950&1,074
\end{tabular}
\end{center}
Operations are live binary additions, subtractions and multiplications;
squares and scalar products are charged, and fixed numeral leaves and
register reuse are free.  The natural-number translations use respectively
441 and 1,081 operations, with unchanged degrees.  These are two different
tradeoffs, rather than a claim that either is globally simplest.

\subsection{The new Pell representation}

Let $\psi_A(0)=0$, $\psi_A(1)=1$ and
$\psi_A(t+2)=2A\psi_A(t+1)-\psi_A(t)$.  Let $\chi_A$ be the companion
sequence with initial values $1,A$.  Thus
\[
 \chi_A(t)^2-(A^2-1)\psi_A(t)^2=1.
\]
For $A\ge2$, the positive Pell solutions are exactly these pairs.

\begin{theorem}[Saturated positive Pell certificate]
\label{nine-saturation}
Let $A,B,C,J$ be positive integers satisfying
$A\ge3$, $1<B<A$, $C>B$ and $J\ge A$.  For $i\ge1$ and $j_0\ge0$ put
\begin{align*}
 D&=1+(A^2-1)C^2,&e_0&=iJC^2,\\
 F&=1+(A^2-1)De_0^2,&\Lambda&=CD+1-A,\\
 G&=A+F\Lambda,&H&=C+F(B+j_0C),\\
 I&=1+(G^2-1)H^2.&&
\end{align*}
Then $C=\psi_A(B)$ if and only if some such $i,j_0$ make $DFI$ a square.
\end{theorem}

\begin{proof}
All displayed radicands and $G,H$ are positive.  In particular
$D\ge A^2$ and $\Lambda\ge A^2-A+1>0$.  The congruences
\begin{align*}
 F&\equiv1\pmod D,&G&\equiv1\pmod D,&I&\equiv1\pmod D,\\
 G&\equiv A\pmod F,&H&\equiv C\pmod F,&I&\equiv D\pmod F
\end{align*}
show that $D,F,I$ are pairwise coprime.  If their product is a square,
each factor is a square.  Set $d=\sqrt D\in\Z_{>0}$.  Then
\[
 F=1+(A^2-1)(iJC^2d)^2.
\]
Pell classification gives positive indices $s,m,t$ such that
\[
 C=\psi_A(s),\qquad iJC^2d=\psi_A(m),\qquad H=\psi_G(t).
\]
The standard rank argument gives $C\mid m$.  Here are its essential
details, including the case of even $C$.  From
$\psi_A(s)\mid\psi_A(m)$ the divisibility law gives $s\mid m$; write
$m=su$.  Expanding the Pell addition identities yields
\[
 \psi_A(su)\equiv
 u\chi_A(s)^{u-1}\psi_A(s)\pmod{\psi_A(s)^3}.
\]
Since $\gcd(\chi_A(s),\psi_A(s))=1$, the stronger divisibility
$\psi_A(s)^2\mid\psi_A(m)$ implies $\psi_A(s)\mid u$, hence $C\mid m$.
The usual divisibility law itself follows by Euclidean reduction of the
addition identities and strict growth of $\psi_A$.

Because $G\equiv A\pmod F$ and $H\equiv C\pmod F$, recurrence reduction
gives
\[
 \psi_A(s)\equiv\psi_A(t)\pmod{\chi_A(m)}.
\]
This implies $s\equiv\pm t\pmod{2m}$.  A short proof avoids any
exceptional-index assumption.  Put $x=\chi_A(m)$.  The addition identities
give
\[
 \psi_A(v+2m)\equiv-\psi_A(v),\qquad
 \psi_A(2m-v)\equiv\psi_A(v)\pmod x.
\]
Every index folds to $r\in[0,m]$, with residue $\pm\psi_A(r)$ and index
congruent to $\pm r\pmod{2m}$.  The Pell identity and $A\ge3$ imply
$x>2\psi_A(m)$.  The signed representatives therefore lie strictly between
$-x/2$ and $x/2$.  Equal residues are equal integers, and strict growth
forces equal folded indices.

On the other hand $G\equiv1\pmod C$ and $H\equiv B\pmod C$, so
$t\equiv B\pmod C$.  Therefore $s\equiv\pm B\pmod C$.  Since $C>1$,
we have $s\ge2$; the inequality $\psi_A(s)\ge As$, together with $B<A$,
gives $0<s+B<C$.  Consequently the only possibility is $s=B$.

Conversely, suppose $C=\psi_A(B)$ and put $d=\chi_A(B)$, so $d^2=D$.
The determinant-one matrix
\[
 M_A=\begin{pmatrix}2A&-1\\1&0\end{pmatrix}
\]
has finite order modulo $JdC^2$.  Choose a positive multiple $m$ of this
order and set $i=\psi_A(m)/(JdC^2)>0$.  Then $F=\chi_A(m)^2$.
Let $T$ be a positive period of $M_G$ modulo $F$.  Choose arbitrarily large
$t\equiv B\pmod{TC}$ and put $H_0=\psi_G(t)$.  We have
\[
 H_0\equiv C\pmod F,\qquad H_0\equiv B\pmod C,
\]
using respectively $G\equiv A\pmod F$ and $G\equiv1\pmod C$.
Since $F\equiv1\pmod C$, the quotient $(H_0-C)/F$ is $B$ modulo $C$.
Taking $t$ sufficiently large makes that quotient at least $B$; hence
\[
 j_0=\frac{(H_0-C)/F-B}{C}\in\N.
\]
Then $H=H_0$ and $I=\chi_G(t)^2$.  This supplies the required witnesses.
All divisions here occur in the witness proof, not in the polynomial.
\end{proof}

The old certificate used
\[
 F_{\rm old}=1+(A^2-1)D^2(i_{\rm old}JC^2)^2.
\]
An old valid witness gives a particularly simple new witness:
$i=\sqrt D\,i_{\rm old}$.  It preserves $F,G,H,I$ exactly.  The converse
is supplied by Theorem~\ref{nine-saturation}, rather than by an unjustified
division of the free integer $i$ by $\sqrt D$.

\subsection{Binding to the actual compact54-sharp source}

Both new equations retain the entire H2082 source prefix and its coherent
compact54-sharp coefficient tuple from commit
\texttt{c1a102936aca437e9999530a0111e6ddab56dc6d}.
The complete layout, three kernel polynomials and coefficient leaves are
included as literal data.  The four original coefficient ports occupy
215,492 raw bytes; the eleven expanded live fixed leaves occupy 1,565,018
raw bytes.  These are different representations of the fixed data and
must not be confused with the number of unknowns.

For the inherited two-block words $S_{\rm code},T_{\rm code}$, define
\begin{align*}
 B_{\rm code}&=2^{864026}b^2,&
 L&=bB_{\rm code}^{2200},\\
 N&=2^{6480184}B_{\rm code}^{1162},&
 R&=L(S_{\rm code}+T_{\rm code}+1)-S_{\rm code}-1,\\
 t_{\rm cap}&=bB_{\rm code}^{116}.&&
\end{align*}
The supplied source tables define both words without a free program index.
The inherited source theorem gives $R,N\ge8$ and the required
carry/binomial calibration.  Define the ratio registers
\begin{align*}
 U&=N^2w,&Y&=N^2s,&M&=RY,\\
 P_0&=2M^2U,&K&=R+1+h(P_0-1),&A&=M(U+1),\\
 B_1&=2R+1,&W_1&=bw,&J&=4A-5,\\
 Q_0&=3W_1C-2(W_1^2-1),&&&\\
 a_0&=1+(P_0^2-1)K^2,&V_0&=P_0K,&&\\
 m_0&=K^2(4t_{\rm cap}-g)-16t_{\rm cap}(C-KY)^2,&&&\\
 \Omega&=Jm_0,&E_0&=J(n+W_1^2)+Q_0.&&
\end{align*}
The nine inputs are
\[
 b,g,h,s,w,C,i,j,n\in\Z_{>0}.
\]
Set $j_0=j-1$ and use $B=B_1$ in Theorem~\ref{nine-saturation}.
The source bounds imply $A\ge33280$, $1<B_1<A$ and $J\ge A$ on every
ambient positive assignment.  The hypothesis $C>B_1$ will be established
from the final equation before the representation theorem is invoked.

\subsection{The 434-operation equation}

Define
\[
 T_0=ADe_0,\qquad V_1=T_0GH,\qquad a_1=DFI,\qquad
 \Theta=E_0-\Omega(V_0+V_1).
\]
The first new equation is
\begin{equation}
 \boxed{\bigl(\Theta^2-\Omega^2(a_0+a_1)\bigr)^2
                    -4\Omega^4a_0a_1=0.}
 \label{nine-cheap-equation}
\end{equation}
Its unconditional majorant is exceptionally small:
\begin{align}
 T_0^2-DF&=D(De_0^2-1)>0,\label{nine-cheap-defect}\\
 V_1^2-a_1&=(T_0^2-DF)(GH)^2+DF(H^2-1)>0.\nonumber
\end{align}
Thus $V_1>\sqrt{a_1}>0$.  Also $V_0^2-a_0=K^2-1>0$.
These inequalities hold before a solution has been assumed.

\begin{proposition}
\label{nine-cheap-equivalence}
Equation~\eqref{nine-cheap-equation} has a positive solution exactly when
the same pinned native source has a solution.
\end{proposition}

\begin{proof}
The identity
\[
 E_0=Jn+(J-2)W_1^2+3W_1C+2>0
\]
shows that a zero cannot have $m_0=0$, since then its left side would be
$E_0^4$.  Factoring the two-root norm over the reals gives signed roots
$r_0^2=a_0$, $r_1^2=a_1$ satisfying
\begin{equation}
 \frac{E_0}{J}=m_0(V_0+V_1+r_0+r_1).
 \label{nine-sign-first}
\end{equation}
The strict majorants make the parenthesis positive for every sign choice;
therefore $m_0>0$ before any square-root integrality is used.

The unchanged margin argument now gives $|C-KY|<K/2$ and
$C>B_1+W_1$.  Moreover $A^2>P_0$, $D>C^2$, $F\ge A^2$ and $I\ge A^2$,
so
\[
 DFI>C^2A^4>K^2P_0^2>a_0.
\]
Hence the rational sum $x=\Theta/\Omega=r_0+r_1$ is nonzero.  The identity
$2xr_0=x^2+a_0-a_1$ makes each root rational, hence integral.
Equation~\eqref{nine-sign-first} then gives $J\mid Q_0$.
Theorem~\ref{nine-saturation} gives $C=\psi_A(B_1)$.
The retained arithmetic ratio/binomial theorem recovers the source zero.
If its original auxiliary certificate is desired explicitly, the finite
period argument in Theorem~\ref{nine-saturation} can be repeated using
the modulus $JDC^2$ instead of $J\sqrt D C^2$.

Conversely, use the inherited canonical ratio witnesses from a source zero.
Set $i=\sqrt D\,i_{\rm old}$ to preserve the old $F,G,H,I$, keeping all
other ratio coordinates fixed.  Choose positive integral roots and put
\[
 n=m_0(V_0+V_1+\sqrt{a_0}+\sqrt{a_1})-W_1^2-Q_0/J.
\]
This is an integer.  The retained bound $W_1^2+Q_0/J<2D$, together with
$V_1>2D$ from $T_0=ADe_0>2D$, proves $n>0$.  One norm factor vanishes.
\end{proof}

Composing this existential equivalence with the named source/refutation
calibration, and then negating, gives the intended ordinary argument in
$\ZFC$.  The new proof does not by itself close the repository's
remaining internal-model calibration or produce a new kernel-checked
object-language consistency biconditional.

\subsection{A smaller-degree majorant of the whole radical}

Temporarily treat $A,C,i,Z$ as independent indeterminates, put $J=4A-5$,
and replace $H$ by $C+FZ$.  Set
\[
 \mathcal Q=DF\bigl((G^2-1)H^2+1\bigr).
\]
Its $A$-degree is 36, with leading coefficient
$(1024C^{19}i^5Z)^2$.  Write $\mathcal Q=\sum q_k A^k$ and define
\begin{align*}
 t_{18}&=1024C^{19}i^5Z,\\
 t_{18-k}&=\frac{q_{36-k}-\sum_{\ell=1}^{k-1}
                  t_{18-\ell}t_{18-k+\ell}}{2t_{18}}
                  \qquad(1\le k\le18).
\end{align*}
This finite calculation takes place in the coefficient field
$\mathbb Q(C,i,Z)$.  It specifies the polynomial part of the formal
square root; it is not a division performed on numerical witnesses.
Exact arithmetic shows that
\[
 \sigma=1024C^5i,\qquad
 \mathcal P=\sigma\sum_{r=0}^{18}t_r A^r
\]
has 333 integer monomials and no negative exponents.  Define
\begin{equation}
 \mathcal L=2^{20}C^{24}i^6Z,\qquad
 \mathcal P_+=\mathcal P+\mathcal L,\qquad
 \mathcal D=\mathcal P_+^2-\sigma^2\mathcal Q.
 \label{nine-radical-polynomial}
\end{equation}
The complete coefficient lists are supplied.  Two exact integer-polynomial
identities give the uniform inequalities:
\begin{itemize}
 \item $\mathcal P(A+3,C+1,i+1,Z+1)$ has 6,650 strictly positive coefficients,
       with constant $15\,915\,642\,730\,652$.
 \item The translated defect $\mathcal D(A+3,C+1,i+1,Z+1)$ has
       36,309 strictly positive coefficients. Its constant is
       $32\,237\,226\,735\,615\,909\,648$.
\end{itemize}
Consequently
\begin{equation}
 \mathcal P_+(A,C,i,Z)>\sigma\sqrt{\mathcal Q}>0
 \quad(A\ge3, C,i,Z\ge1).
 \label{nine-whole-bound}
\end{equation}
This is a universal positivity certificate.  The uncorrected formal
polynomial part would be invalid: its defect at $(3,1,1,1)$ is
$-1\,140\,296\,347\,768\,021\,232$.  The correction in
\eqref{nine-radical-polynomial} is therefore substantive.

For the second fixed equation put
\[
 Z=B_1+(j-1)C,\quad
 V_1=\mathcal P_+(A,C,i,Z),\quad a_1=(1024C^5i)^2DFI,
\]
and use exactly the norm equation~\eqref{nine-cheap-equation}.
The previous proof applies with one additional observation: once the
signed root $r_1$ is integral, $r_1/\sigma$ is a rational square root of
the integer $DFI$, hence is itself an integer.  Thus the scale introduces
no extra square solutions.  In the converse direction $F>D$ and
$\sigma\ge1024$ imply $V_1>2D$, so the witness $n$ remains positive.

The supplied graph evaluates the 333 coefficients by fixed polynomial
operations.  All 24 multivariate Horner orders were counted; the selected
order is $Z,i,C,A$.  This yields 1,074 live binary operations.  It is a
finite evaluated circuit count, not a claim of an optimal arithmetic circuit.

\subsection{Exact degrees and the changed dominant term}

The source has
\[
 r=\deg R=8799,\qquad \nu=\deg N=2324,\qquad W=116,
 \qquad a=r+4\nu+2=18097.
\]
The common register degrees are
\begin{align*}
 \deg\Omega&=5a+2W+3=90720,\\
 \deg V_0&=3a+r+1=63091,\\
 \deg\delta_0&=3a+r+2=63092,\qquad \delta_0=K^2-1,\\
 \deg E_0&=a+4=18101,\qquad\deg a_0=126182.
\end{align*}
For either majorant write $\delta_1=V_1^2-a_1$ and
\[
 B_* =E_0^2-2E_0\Omega(V_0+V_1)+\Omega^2(\delta_0+\delta_1).
\]
The universal exact identity
\begin{equation}
 F_9=B_*^2+4\Omega^2V_0V_1B_*
           +4\Omega^4(\delta_0V_1^2+a_0\delta_1)
 \label{nine-resolved-norm}
\end{equation}
resolves the leading cancellations that a raw syntactic degree count misses.

For the simple majorant,
\[
 \deg V_1=18a+r+24=334569,
 \qquad\deg\delta_1=34a+2r+48=632944.
\]
The latter follows directly from~\eqref{nine-cheap-defect}.  The square
$B_*^2$ is uniquely dominant in~\eqref{nine-resolved-norm}; therefore
\[
 \deg F_9=92r+352\nu+8W+284=1\,628\,768.
\]
Its leading coefficient is positive on the actual source tuple.

For the complete-radical bound the unique greatest weighted monomials,
with weights $(a,1,1,r)$ on $(A,C,i,Z)$, are
\[
 2^{20}A^{18}C^{24}i^6Z\quad\hbox{in }\mathcal P_+,
 \qquad
 2^{41}A^{18}C^{48}i^{12}Z^2\quad\hbox{in }\mathcal D.
\]
Hence $\deg V_1=334575$, $\deg\delta_1=343404$ and
$\deg B_*=524844$.  The four summands in~\eqref{nine-resolved-norm},
splitting its last parenthesis, have degrees
\begin{center}
\begin{tabular}{lr}
$B_*^2$&1,049,688\\
$4\Omega^2V_0V_1B_*$&1,103,950\\
$4\Omega^4\delta_0V_1^2$&1,095,122\\
$4\Omega^4a_0\delta_1$&832,466
\end{tabular}
\end{center}
The second term is now uniquely dominant, with positive leading
coefficient.  Thus the exact degree is
\[
 \boxed{\deg F_9=63r+236\nu+8W+221=1\,103\,950.}
\]
Translation of positive inputs to natural successors preserves these
degrees; the natural $j$ directly supplies the existing $j-1$ coordinate.

\subsection{What the executable evidence proves}

The independent verifier uses a separately implemented sparse integer
polynomial ring and does not import either circuit builder or the rational
discovery code.  It checks the entire radical identity, shifted positive
coefficient lists, all local certificate formulas, the exact final norm,
source-prefix structural equality, topology, liveness, operation counts,
all fixed-leaf hashes and the resolved degree comparisons.  It rejects
both the uncorrected radical truncation and a corrupted final norm.
The complete coefficient generators and parent proofs are retained as
explicit dependencies; no finite check is presented as a consistency test.

The main mathematical change is the saturation theorem, whose soundness
and finite-period completeness were proved above.  The whole-radical
certificate is a second new uniform lemma.  The remaining native/compiler
and syntactic-refutation interfaces are inherited, and their unresolved
formal verification boundary remains visible.
